\section{Resonant Network and Transformer Ratio}
\label{sec:analysis-resonant}

The output-tank calculation adds the primary-referred series contributions: approximately \SI{9.4}{\micro\henry} from the discrete resonant inductor and approximately \SI{16.5}{\micro\henry} from transformer leakage. Together they give
\begin{equation}
\LRES\approx\SI{25.9}{\micro\henry}.
\end{equation}
With \(\CRES=\SI{1}{\nano\farad}\), the nominal series resonance is
\begin{equation}
f_0\approx\frac{1}{2\pi\sqrt{\LRES\CRES}}
\approx\SI{0.989}{\mega\hertz}.
\end{equation}
The nominal resonance lies close to the \SI{1}{\mega\hertz} carrier. Loaded gain and commutation also depend on the magnetizing branch and pump current.

The transformer turns ratio is defined as
\begin{equation}
n=\frac{N_p}{N_s},
\end{equation}
using the same convention as Equation~\eqref{eq:app-turns-ratio}. The supplied transformer measurements are near 8:1, so \(n\approx8\) is used for the step-down ratio. The reference LED-driver equations use \(N_s/N_p\), which must be inverted before applying their relationships here.

With full-wave rectification and approximately sinusoidal current, the equivalent load at the secondary rectifier input is
\begin{equation}
R_{\mathrm{REC,s}}=\frac{8V_{\mathrm{sec}}^2}{\pi^2P_{\mathrm{sec}}},
\end{equation}
Referring this resistance to the primary gives
\begin{equation}
R_{\mathrm{REC,p}}=n^2R_{\mathrm{REC,s}}.
\end{equation}
Here \(P_{\mathrm{sec}}\) includes downstream demand and losses. Using connector power alone would understate tank loading.

If magnetizing current is neglected, the required series-tank gain is
\begin{equation}
M_{\mathrm{RES}}=\frac{2nV_{\mathrm{sec}}}{\VDC}.
\end{equation}
Using \(f_n=f_s/f_0\) and \(Q_L=\omega_0\LRES/R_{\mathrm{REC,p}}\), the reference first-harmonic gain is
\begin{equation}
M_{\mathrm{RES}}=
\frac{1}{\sqrt{1+Q_L^2(f_n-f_n^{-1})^2}}.
\label{eq:analysis-series-gain}
\end{equation}
The associated component relationships are
\begin{equation}
\LRES=\frac{Q_LR_{\mathrm{REC,p}}}{\omega_0},
\qquad
\CRES=\frac{1}{\omega_0Q_LR_{\mathrm{REC,p}}}.
\end{equation}
These first-harmonic relationships provide an initial component and gain check at the assumed load \cite{ammar2021led}.

The supplied magnetizing inductance is approximately \SIrange{140}{151}{\micro\henry}, about five to six times nominal series inductance. Its effect on input impedance increases at light load and burst transitions. Because Equation~\eqref{eq:analysis-series-gain} omits this branch, measured gain discrepancies require a fuller model before revising series components.

Installed behavior also depends on transformer variant and polarity, rectifier conduction, and frequency-dependent parasitics. Chapters~\ref{ch:magnetics} and \ref{ch:hardware} describe the integration.
